
\documentclass[tikz,border=2pt]{standalone}
\usepackage[utf8]{inputenc}
\usepackage[T1]{fontenc}
\usepackage[french]{babel}
\usepackage{amsmath,amssymb}
\usepackage[table]{xcolor}
\usepackage{siunitx}
\usepackage[version=4]{mhchem}
\usepackage{tabularx,booktabs,array}
\usepackage{tikz}
\usetikzlibrary{arrows.meta,positioning,calc,fit,decorations.pathmorphing,patterns,shapes.geometric}
\usepackage{pgfplots}
\pgfplotsset{compat=1.18}
\definecolor{primary}{RGB}{14,116,144}
\definecolor{accent}{RGB}{16,185,129}
\definecolor{dark}{RGB}{15,23,42}
\definecolor{bglight}{RGB}{241,245,249}
\definecolor{borderlight}{RGB}{203,213,225}
\definecolor{examplepink}{RGB}{236,72,153}
\definecolor{remarkblue}{RGB}{14,165,233}
\definecolor{notepurple}{RGB}{99,102,241}
\definecolor{synthgreen}{RGB}{5,150,105}
\definecolor{synthorange}{RGB}{249,115,22}
\definecolor{synthviolet}{RGB}{79,70,229}
\definecolor{primary}{RGB}{14,116,144}
\definecolor{accent}{RGB}{16,185,129}
\definecolor{dark}{RGB}{15,23,42}
\definecolor{bglight}{RGB}{241,245,249}
\definecolor{borderlight}{RGB}{203,213,225}
\definecolor{examplepink}{RGB}{236,72,153}
\definecolor{remarkblue}{RGB}{14,165,233}
\definecolor{notepurple}{RGB}{99,102,241}
\definecolor{synthgreen}{RGB}{5,150,105}
\definecolor{synthorange}{RGB}{249,115,22}
\definecolor{synthviolet}{RGB}{79,70,229}
\definecolor{primary}{RGB}{14,116,144}
\definecolor{accent}{RGB}{16,185,129}
\definecolor{dark}{RGB}{15,23,42}
\definecolor{bglight}{RGB}{241,245,249}
\definecolor{examplepink}{RGB}{236,72,153}
\definecolor{remarkblue}{RGB}{14,165,233}
\definecolor{notepurple}{RGB}{99,102,241}
\definecolor{synthgreen}{RGB}{5,150,105}
\definecolor{synthorange}{RGB}{249,115,22}
\definecolor{synthviolet}{RGB}{79,70,229}
\begin{document}
\begin{tikzpicture}[>=Stealth, scale=0.9]
 % Photocathode
 \draw[thick, fill=bglight] (0,0) rectangle (0.4,3);
 \node[rotate=90] at (0.2,1.5) {\small Photocathode};

 % Dynodes
 \foreach \i/\y in {1/0.5, 2/1.2, 3/1.9, 4/2.6} {
 \draw[thick, fill=borderlight] (1.2+\i*1.2, \y-0.15) rectangle (1.2+\i*1.2+0.3, \y+0.8);
 \node[font=\tiny] at (1.35+\i*1.2, \y+0.32) {D\textsubscript{\i}};
 }

 % Anode
 \draw[thick, fill=primary!20] (7.8,1) rectangle (8.2,2.2);
 \node[rotate=90, font=\small] at (8.0,1.6) {Anode};

 % Photon entrant
 \draw[->, very thick, color=synthorange, line width=1.5pt] (-1.5,1.5) -- (-0.1,1.5);
 \node[color=synthorange, font=\small] at (-1.5,1.9) {Photon};
 \node[color=synthorange, font=\small] at (-1.5,1.1) {incident};

 % Electrons (schématique)
 \draw[->, thick, dashed, color=examplepink] (0.5,1.5) -- (2.2,0.7);
 \draw[->, thick, dashed, color=examplepink] (2.7,0.9) -- (3.3,1.5);
 \draw[->, thick, dashed, color=examplepink] (3.3,1.5) -- (3.5,1.6);
 \draw[->, thick, dashed, color=examplepink] (3.8,1.6) -- (4.5,2.2);
 \draw[->, thick, dashed, color=examplepink] (4.5,2.2) -- (4.7,2.3);
 \draw[->, thick, dashed, color=examplepink] (5.0,2.5) -- (5.8,2.8);
 \draw[->, thick, dashed, color=examplepink] (6.2,2.9) -- (7.7,1.6);

 \node[color=examplepink, font=\tiny] at (4.5,0) {\'{e}lectrons secondaires};
\end{tikzpicture}
\end{document}
